\begin{table}[ht]
\centering
\caption{Litigation-specificity placebo battery}
\label{tab:h7_placebo_battery}
\begin{threeparttable}
\small
\setlength{\tabcolsep}{4.5pt}
\begin{tabular}{lccccc}
\toprule
 & All never-lit. & Matched never-lit. & Matched + qty & Matched ref. price & Matched bidders \\
\midrule
Urgent purchase & -0.020 & -0.046 & -0.020 & -0.071 & -0.032 \\
 & (0.032) & (0.053) & (0.070) & (0.056) & (0.043) \\
\addlinespace
Observations & 39{,}283 & 3{,}915 & 3{,}915 & 5{,}215 & 2{,}490 \\
Item FE & Yes & Yes & Yes & Yes & Yes \\
Year FE & Yes & Yes & Yes & Yes & Yes \\
PBU FE & Yes & Yes & Yes & Yes & Yes \\
Quantity control & No & No & Yes & No & No \\
\bottomrule
\end{tabular}
\begin{tablenotes}[flushleft]\footnotesize
\item \textit{Notes:} The matched placebo restricts never-litigated items to buyer-by-class environments that also contain litigated purchases elsewhere in the panel. The exercise is a specificity check: it asks whether the urgent-procurement pattern appears where the litigation channel is absent but procurement environments overlap. Price columns use accepted winning bids when the outcome is negotiated price. Standard errors are clustered by PBU. $^{*}p<0.10$, $^{**}p<0.05$, $^{***}p<0.01$.
\end{tablenotes}
\end{threeparttable}
\end{table}
